% Figure: a finite/infinite square-well sketch showing the confined
% conduction-subband levels of a quantum well. The well of width L holds
% discrete levels E1, E2, E3 that scale as n^2/L^2; the emission photon
% energy is the gap plus the electron and hole confinement energies, so
% it shifts blue as the well narrows. Plain TikZ.
\begin{figure}[htbp]
\centering
\begin{tikzpicture}[font=\scriptsize]
% well walls (conduction band): barrier high on both sides, well in middle
\draw[engblue, very thick] (0,3.4) -- (2,3.4);        % left barrier top
\draw[engblue, very thick] (2,3.4) -- (2,0.6);        % left wall down
\draw[engblue, very thick] (2,0.6) -- (5,0.6);        % well floor
\draw[engblue, very thick] (5,0.6) -- (5,3.4);        % right wall up
\draw[engblue, very thick] (5,3.4) -- (7,3.4);        % right barrier top
% well-width marker
\draw[<->, enggray] (2,0.35) -- (5,0.35);
\node[enggray, anchor=north, font=\tiny] at (3.5,0.32) {well width $L$};
% confined levels inside well: E1, E2, E3 (n^2 spacing)
\draw[engred, thick] (2.1,1.0) -- (4.9,1.0);
\node[engred, anchor=west, font=\tiny] at (5.05,1.0) {$E_1 \propto 1/L^2$};
\draw[engred, thick] (2.1,1.9) -- (4.9,1.9);
\node[engred, anchor=west, font=\tiny] at (5.05,1.9) {$E_2 = 4E_1$};
\draw[engred, thick] (2.1,3.1) -- (4.9,3.1);
\node[engred, anchor=west, font=\tiny] at (5.05,3.1) {$E_3 = 9E_1$};
% ground-state half-wave probability sketch
\draw[engorange, thick, domain=2:5, samples=40, smooth]
  plot (\x, {1.0 + 0.55*sin(deg((\x-2)*pi/3))});
\node[engorange, anchor=south, font=\tiny] at (3.5,1.55) {$|\psi_1|$};
% band-gap / emission arrow on the left
\node[anchor=east, font=\tiny, align=right] at (1.9,3.4) {barrier};
\node[anchor=east, font=\tiny, align=right] at (1.9,0.6) {$E_c$ (bulk)};
\draw[->, engamber, thick] (1.3,3.2) -- (1.3,0.8);
\node[engamber, anchor=east, font=\tiny, align=right] at (1.2,2.0)
  {$\hbar\omega = E_g + E_1^e + E_1^h$};
\end{tikzpicture}
\caption[Quantum-well confined levels and emission]{A quantum well is a thin
layer of narrow-gap material between wider-gap barriers, confining the
electron in one direction. The confinement quantises the motion across the
well into discrete subband levels $E_n \propto n^2/L^2$, the particle-in-a-box
result. The emitted photon energy is the bulk gap plus the electron and hole
ground-state confinement energies, so narrowing the well shifts the emission
to the blue. Choosing $L$ tunes the laser or LED wavelength at fixed material,
and concentrating the available states near the band edge lowers the laser
threshold current, the design payoff noted in the text.}
\label{fig:vol04ch12:quantum-well-levels}
\end{figure}
